% vim:foldmethod=marker:foldmarker=(((,))):foldlevel=0

\subimport{.preamble}{packages.tex}

%: External variables (((

% russian names for autoref:
\AtBeginDocument{\renewcommand\figureautorefname{рис.}}
\AtBeginDocument{\renewcommand\theoremautorefname{теор.}}
\AtBeginDocument{\renewcommand\lemmaautorefname{лемма}}
\AtBeginDocument{\renewcommand\definitionautorefname{опр.}}
\AtBeginDocument{\renewcommand\statementautorefname{утв.}}

% BibTeX month names:
% \def\bbljan{Jan}
% \def\bblfeb{Feb}
% \def\bblmar{Mar}
% \def\bblapr{Apr}
% \def\bblmay{May}
% \def\bbljun{Jun}
% \def\bbljul{Jul}
% \def\bblaug{Aug}
% \def\bblsep{Sep}
% \def\bbloct{Oct}
% \def\bblnov{Nov}
% \def\bbldec{Dec}

%: )))

%: Math operators (((

% constant
\DeclareMathOperator{\const}{const}

% matrix rank
\DeclareMathOperator{\rk}{rk}

% image
\DeclareMathOperator{\Img}{Im}

% sign
\DeclareMathOperator{\sign}{sign}

% Divergence
\DeclareMathOperator{\dive}{div}

% hyperbolic functions
\DeclareMathOperator{\ch}{ch}
\DeclareMathOperator{\sh}{sh}

% differential
\newcommand\di[1][]{\ensuremath{\operatorname{d}^{#1}}}

% strong limit
\DeclareMathOperator{\slim}{s-lim}

%: )))

%: Equations and labels (((

% labeled equation
\newenvironment{equ}[1]{\begin{equation} \label{eq:#1}}{\end{equation}}
% multiline equation with *every* line numbered
\newenvironment{mequ}[1][\empheqlbrace]
{
    \setkeys{EmphEqEnv}{align} \setkeys{EmphEqOpt}{left={#1}} \EmphEqMainEnv
}
{\endEmphEqMainEnv}

% labels and references
\newcommand*\lbl[1]{\label{eq:#1}}
\newcommand*\eref[1]{\eqref{eq:#1}}

%: )))

%: Shorthands (((

\newcommand\sub{\ensuremath\subset}

% phi and epsilon
\let\tphi\phi
\let\phi\varphi
\let\teps\epsilon
\let\eps\varepsilon

% short for stmaryrd's \llbracket and \rrbracket
\newcommand\llb\llbracket
\newcommand\rrb\rrbracket

% over-decorations
\let\ol\overline
\let\vawe\widetilde
% "raised overline" -- for capital letters
\newcommand*{\roverline}[2][1pt]{%
	% Namespace for \argheight
	\begingroup%
	\newlength\argheight%
	\settoheight{\argheight}{#2}%
	\addtolength{\argheight}{#1}%
	\ensuremath{\overline{\raisebox{0pt}[\argheight]{\ensuremath{{#2}}}}}%
	\endgroup%
}
% long hat
\newcommand\smallhat\hat
\AtBeginDocument{\renewcommand\hat\widehat}

\newcommand\op\operatorname

% limits
\newcommand\ulim{\ensuremath{\varlimsup\limits}}
\newcommand\dlim{\ensuremath{\varliminf\limits}}

% := (DEPRECATED: use /coloneq)
\newcommand\define\coloneqq
% =: (DEPRECATED: use \eqqcolon)
\newcommand\fed\eqqcolon

% set package shorthands (DEPRECATED: prefer Set and Braket)
\let\tset\set
\let\set\Set

%: )))

%: Bells and whistles (and misc) (((

% imaginary unit
\newcommand\ii{\mathfrak i}

% text in a circle
\newcommand\circled[1]{\text{\textcircled{$ {#1} $}}}

% partial derivative
\newcommand*\pder[2][]{\ensuremath{\frac{\partial {#1}}{\partial {#2}}}}

% n-th derivative
\newcommand*\nder[2][n]{\ensuremath{{#2}^{({#1})}}}

% variants
\newenvironment{vars}{\left[\begin{matrix}}{\end{matrix}\right.}

% by definition (default is =)
\newcommand*\bydef[1][=]{\ensuremath{\stackrel{\operatorname{def}}{#1}}}

% column
\newcommand\column[3][bmatrix]{\ensuremath{
        \begin{#1}
            {#2} \\
            \vdots \\
            {#3}
        \end{#1}
    }}

% divisible by
\DeclareRobustCommand{\divby}{%
  \mathrel{\vbox{\baselineskip.65ex\lineskiplimit0pt\hbox{.}\hbox{.}\hbox{.}}}%
}

% not divisible by
\DeclareRobustCommand{\ndivby}{%
  \mathrel{\not\vbox{\baselineskip.65ex\lineskiplimit0pt\hbox{.}\hbox{.}\hbox{.}}}%
}

% controversy
\newcommand\contra{~\scalebox{1.2}{$\lightning$}~}

% roman numerals
\makeatletter
\newcommand*{\rom}[1]{\text{\expandafter\@slowromancap\romannumeral #1@}}
\makeatother

% triangle inequality
\newcommand\trile{\ensuremath{\stackrel\vartriangle\le}}
\newcommand\trige{\ensuremath{\stackrel\vartriangle\ge}}

% math &
\newcommand\amp{\ensuremath{\mathrel\&}}

%: )))

%: Extensible arrows (((

% arrows with text below `em
\newcommand\underarr[1]{\ensuremath\xrightarrow[#1]{}}
\newcommand\underimp[1]{\ensuremath\xRightarrow[#1]{}}
\newcommand\underiff[1]{\ensuremath\xLeftrightarrow[#1]{}}

% arrows with text above `em (DEPRECATED: it's pointless)
\newcommand\overarr[1]{\ensuremath\xrightarrow{#1}}

% arrow-like with text below `em
\newcommand\undereq[1]{\ensuremath\xlongequal[#1]{}}

\makeatletter
\let\equiv@\equiv
\newcommand\@equiv[1][]{\overset{#1}\equiv@}
\let\equiv\@equiv
\makeatother

% custom extensible arrows
\makeatletter
% xrightrightarrow
\newdimen\@temp
\settowidth\@temp{$ \m@th \relbar $}
\def\@highRelbar{\raisebox{-0.1675em}{$ \relbar $}\hspace{-\@temp}\raisebox{0.1675em}{$ \relbar $}}
\def\xrightrightarrowsfill@{\arrowfill@\@highRelbar\@highRelbar\rightrightarrows}
\newcommand\xrightrightarrows[2][]{\ext@arrow 0099\xrightrightarrowsfill@{#1}{#2}}

% xequiv
\let\equv@\equiv
\newcommand\xequivfill@{\arrowfill@\equiv@\equiv@\equiv@}
\newcommand\xequiv[2][]{\ext@arrow 0099\xequivfill@{#1}{#2}}
% equiv with domain
\AtBeginDocument{\renewcommand\equiv[1][]{\xequiv{#1}}}

\makeatother

% special cases
\newcommand\uniarr[2][n \to \infty]{\ensuremath\xrightrightarrows[#1]{#2}}

% other custom arrows

\newcommand\curveddownarrowright{\rotatebox{180}{$ \curvearrowleft $}}
\newcommand\curveddownarrowleft{\rotatebox{180}{$ \curvearrowright $}}

% 1 is left-to-right
\newcommand\curvedir[2][1]{\ifthenelse{\isodd{#1}}
    {\ensuremath{\overset{\curveddownarrowright}{#2}}}
    {\ensuremath{\overset{\curveddownarrowleft}{#2}}}}

%: )))

%: Sets (((

\newcommand\N{\ensuremath{\mathbb{N}}}
\newcommand\R{\ensuremath{\mathbb{R}}}
\newcommand\Z{\ensuremath{\mathbb{Z}}}
\renewcommand\O{\ensuremath{\emptyset}}
\newcommand\Q{\ensuremath{\mathbb{Q}}}
\newcommand\Co{\ensuremath{\mathbb{C}}}     % because /C is already taken
\newcommand\Prime{\ensuremath{\mathbb{P}}}     % prime numbers
\newcommand\On[1][]{\ensuremath{\mathbb{O}_{#1}}}  % (0, ..., 0)

%: )))

%: Text abbreviations (((

\newcommand\as{\text{т.~к.}\xspace}
\newcommand\As{\text{Т.~к.}\xspace}
\newcommand\ie{\text{т.~е.}\xspace}
% so-called
\newcommand\soc{\text{т.~н.}\xspace}
% always-everywhere
\newcommand\ale{\text{п.~в.}\xspace}

%: )))

%: Source-related (((

\makeatletter
\newcommand\TODO[1]{\textcolor{red}{{\Huge TODO:} {#1}}
\@latex@warning{TODO: #1}
}
\makeatother

\newcommand\comment[1]{\textcolor{magenta}{\textit{#1}}}

%: )))

\subimport{.preamble}{theorems.tex}

\subimport{..}{hyphenations.tex}
